\chapter*{Conclusion}

As we saw, this survey analyzed the role of various AI tools in supporting software development, mainly focusing on their impact on software quality. By examining various categories, such as code generating, analysis, and test automation, the work aimed to give an overview of how these technologies are being used in today's software engineering practices.

The review highlighted how AI can be a positively used tool in various fields, such as maintainability, reliability, and bug prevention. Code generating tools help the developers' productivity, while code analysis and test automation ones have a more important role in finding bugs, making sure company standards are enforced, and improving the tests effectiveness. It is important to remember, however, that such benefits usually come with multiple trade-offs related to control, transparency, and long-term maintainability.

A theme that could've been found all throughout the survey is how much human oversight is needed. Despite AI tools being able to act autonomously and enhance various steps of the development cycle, their results and recommendations need a careful review by human before being integrated in the main code structure, especially to avoid over-reliance and quality degradation. Because of this, concepts like technical \textit{debt} and \textit{trust} are important  when evaluating the long-term results of using AI in the development cycle

Overall, AI tools should be used in a way that helps programmers and does not replace them. When used properly in workflows and reviewed under the appropriate criteria, they have the potential to make software more reliable and maintainable. A continued research is going to be key for a better understanding on the impact they yield long-term on software quality and development practices.